% 中英文摘要
\phantomsection
\pagenumbering{Roman}
\addcontentsline{toc}{chapter}{摘\hspace{1em}要}

\begin{cnabstract}{预期共识效应；资产定价；Agent；强化学习}

金融市场是一个多方投资者参与的市场，具有复杂性、动态性等特点。
投资者在单个资产上形成一致的预期，能否影响资产价格，
是一个理解金融市场的运行机制和金融市场的稳定的重要问题。
本文中构建了多个Agent投资者，使用88个因子数据作为输入，MLP模型作为Agent的判断规则，
由强化学习训练Agent，试图探究Agent投资者的预期共识定价效应。
实验的结果显示投资者的预期共识度越高，股票的未来收益越大，
这一结论在一系列稳健性检验后依然成立。
进一步的，本文从融资和融券约束的角度研究预期共识定价效应的形成机制。
实验的结果表明了融资和融券约束是预期共识定价效应的重要形成机制，
在放松融资、融券限制后，预期共识定价效应变弱。
本文还从企业市值、股权异质性和市场周期的角度，进一步研究了预期共识定价效应的异质性。
结果表明，小市值企业、非国企企业和牛市周期中的企业，其预期共识定价效应更高。
最后，本文就研究结果提出相应的政策建议和研究展望。

  
\end{cnabstract}
\newpage

\phantomsection
\addcontentsline{toc}{chapter}{ABSTRACT}

\begin{enabstract}{Expectation consensus effect; Asset pricing; Agent; Reinforcement learning}

With the rapid development of artificial intelligence, investors can process information and make trading decisions more precisely. 
Whether investors' expectation consensus on specific assets affects asset prices 
is critical to understanding the operation and stability of financial markets.
This paper constructs multi-agent investors with 88 factors as input, 
adopts the MLP model as the agents' decision-making framework, 
and explores the pricing effect of agents' expectation consensus through reinforcement learning training.
Empirical results show that higher investor expectation consensus corresponds 
to higher future stock returns, a conclusion that remains robust after multiple checks. 
Further analysis finds that margin trading and short-selling constraints are the core formation mechanism of this pricing effect, 
which weakens significantly when such constraints are relaxed. We also document significant heterogeneity: 
the effect is far more pronounced for small-cap firms, non-state-owned enterprises, and firms in bull market cycles. 
Finally, we propose corresponding policy recommendations and research prospects based on the findings.
\end{enabstract}
\newpage
